% 攻撃者モデル図 (Threat Model Diagram)
% 独立可編集ファイル。lualatex で出力。
\documentclass[border=5mm]{standalone}

\usepackage{luatexja}
\usepackage{tikz}
\usetikzlibrary{positioning, arrows.meta}

\begin{document}

\begin{tikzpicture}[
    box/.style={draw, rounded corners, align=center,
                minimum width=26mm, minimum height=12mm, font=\small},
    attacker/.style={draw, rounded corners, align=center,
                minimum width=32mm, minimum height=18mm,
                fill=red!8, font=\small},
    channel/.style={draw, dashed, rounded corners, align=center,
                minimum width=22mm, minimum height=10mm,
                fill=blue!5, font=\footnotesize},
    small/.style={font=\footnotesize},
    tiny/.style={font=\scriptsize},
    >=Latex
  ]

  % ===== 主經路：送信機 → 無線チャネル → 受信機 =====
  \node[box] (TX) {送信機\\(Controller)};
  \node[channel, right=14mm of TX] (CH) {無線チャネル};
  \node[box, right=14mm of CH] (RX) {受信機\\(Device)};

  % 正規通信の矢印
  \draw[-Latex, thick] (TX) -- (CH);
  \draw[-Latex, thick] (CH) -- (RX);

  % ===== 攻撃者（チャネルの下に配置） =====
  \node[attacker, below=20mm of CH] (A)
    {攻撃者 A\\[2pt]
     {\scriptsize Listen / Record / Replay}};

  % 傍受の破線矢印（チャネル → 攻撃者）
  \draw[dashed, -Latex, thick, red!60!black]
    (CH.south) -- node[small, left, pos=0.5]{傍受} (A.north);

  % 再送の破線矢印（攻撃者 → 受信機）
  \draw[dashed, -Latex, thick, red!60!black]
    (A.east) to[bend left=25]
    node[small, right, pos=0.5]{再送} (RX.south);

  % ===== 図下の注記 =====
  \node[tiny, below=6mm of A]
    {想定：鍵は不所持だが，傍受・再送は可能};

\end{tikzpicture}

\end{document}
